\chapter{Tujuan dan Arsitektur Sistem}

\section{Tujuan Pengembangan}

Tujuan utama pengembangan AuditChain Gateway adalah sebagai berikut:

\begin{enumerate}
    \item Menjamin integritas data pada sistem klien tanpa mengubah kode aplikasi existing secara langsung atau zero code intrusion pada aplikasi utama.
    \item Menyediakan audit trail yang tidak dapat dimanipulasi secara diam-diam atau tamper-evident menggunakan kombinasi hashing kriptografis SHA3-256, Merkle Tree, dan blockchain anchoring.
    \item Memberikan dashboard monitoring dan verifikasi bagi auditor, admin, dan pihak yang bertanggung jawab terhadap integritas data.
    \item Mendukung berbagai jenis basis data klien, termasuk PostgreSQL, MySQL, Oracle, MongoDB, dan SQL Server.
\end{enumerate}

\section{Komponen Arsitektur}

Sistem terdiri dari beberapa komponen utama:

\begin{itemize}
    \item \textbf{Ingestion Layer} menerima data perubahan INSERT, UPDATE, dan DELETE dari database klien melalui Kafka CDC berbasis Debezium.
    \item \textbf{Hasher Engine} menghitung hash SHA3-256 dari setiap log yang masuk.
    \item \textbf{Aggregator Engine} mengelompokkan log menjadi batch dan membangun Merkle Tree.
    \item \textbf{Fabric Anchoring} mengirimkan Merkle Root ke jaringan Hyperledger Fabric.
    \item \textbf{Verification Engine} melakukan verifikasi integritas secara berlapis, mulai dari database existence, local re-hash, agent source verification, sampai blockchain consensus.
    \item \textbf{Dashboard Web} berbasis React menyediakan antarmuka bagi admin dan auditor untuk memonitor, mencari, memfilter, dan memverifikasi data.
\end{itemize}

\dataflowdiagram

\section{Lifecycle Status Log}

Setiap log yang masuk ke sistem melewati beberapa status utama. Log pertama kali diterima sebagai event, kemudian dihitung hash-nya, dikelompokkan dalam batch Merkle Tree, dan setelah itu Merkle Root dikirim ke blockchain.

\lifecyclediagram

Lifecycle ini penting karena dashboard dapat membedakan log yang baru diterima, log yang sudah di-hash, log yang sudah masuk batch, log yang sudah di-anchor, serta log yang perlu diperiksa karena statusnya tampered, pending, atau unreachable.
